\documentclass[11pt, letterpaper]{article}

% ── Packages ──────────────────────────────────────────────────────────────────
\usepackage[margin=1in]{geometry}
\usepackage{amsmath, amssymb}
\usepackage{booktabs}
\usepackage{graphicx}
\usepackage{hyperref}
\usepackage{xcolor}
\usepackage{enumitem}
\usepackage{microtype}
\usepackage{parskip}
\usepackage{titlesec}
\usepackage{fancyhdr}
\usepackage{caption}
\usepackage{array}
\usepackage{multirow}

% ── Colors ────────────────────────────────────────────────────────────────────
\definecolor{lewisblue}{RGB}{0, 56, 101}
\definecolor{darkgray}{RGB}{60, 60, 60}
\definecolor{lightgray}{RGB}{240, 240, 240}

% ── Hyperref ──────────────────────────────────────────────────────────────────
\hypersetup{
  colorlinks=true,
  linkcolor=lewisblue,
  urlcolor=lewisblue,
  citecolor=lewisblue,
  pdftitle={Trustworthy Real-Time Air Quality Monitoring Using Low-Cost PM2.5 Sensor Networks},
  pdfauthor={Idhant Ranjan, Arun Muthukumar, Varun Kalidindi}
}

% ── Section formatting ────────────────────────────────────────────────────────
\titleformat{\section}{\large\bfseries\color{lewisblue}}{}{0em}{}[\titlerule]
\titleformat{\subsection}{\normalsize\bfseries\color{darkgray}}{}{0em}{}
\titleformat{\subsubsection}{\normalsize\itshape\color{darkgray}}{}{0em}{}

% ── Header / Footer ───────────────────────────────────────────────────────────
\pagestyle{fancy}
\fancyhf{}
\renewcommand{\headrulewidth}{0.4pt}
\fancyhead[L]{\small\color{darkgray}ACS Air Quality Project --- Lewis University}
\fancyhead[R]{\small\color{darkgray}Spring 2026}
\fancyfoot[C]{\small\thepage}

% ── Title ─────────────────────────────────────────────────────────────────────
\title{%
  \vspace{-1cm}
  {\Large \textbf{Trustworthy Real-Time Air Quality Monitoring}}\\[4pt]
  {\large Using Low-Cost PM$_{2.5}$ Sensor Networks, Machine Learning,\\
  and Calibrated Uncertainty Estimation}
}
\author{%
  Idhant Ranjan \quad Arun Muthukumar \quad Varun Kalidindi\\[4pt]
  \small Lewis University\\[2pt]
  \small Mentored by Dr.\ Kozminski
}
\date{Spring 2026}

% ──────────────────────────────────────────────────────────────────────────────
\begin{document}
\maketitle
\thispagestyle{fancy}

% ── Abstract ──────────────────────────────────────────────────────────────────
\begin{abstract}
Low-cost PM$_{2.5}$ sensors provide high-frequency, community-scale air quality data but
introduce reliability challenges due to calibration drift and hardware variability.
We present an end-to-end air quality monitoring system built on a network of five
dual-channel sensors deployed at Lewis University during November 2025.
Our contributions are threefold: (1) a \textbf{dual-sensor trust scoring} framework
that quantifies reading reliability from the A/B channel agreement of each sensor unit;
(2) a \textbf{rule-based event detection} pipeline that identifies smoke spikes,
precipitation washout events, and sensor failures from 2-minute-resolution time series;
and (3) a \textbf{Random Forest forecasting model} with tree-bootstrap uncertainty
intervals (80\% coverage) that predicts PM$_{2.5}$ one step ahead.
The forecasting model achieves a mean absolute error of 0.51--1.12~$\mu$g/m$^3$ and
$R^2$ of 0.77--0.98 across the four high-coverage sensors.
Results are surfaced in an interactive React dashboard that includes per-sensor
trust badges, EPA AQI color coding, an event feed, and short-term forecast charts
with calibrated confidence bands.
\end{abstract}

\tableofcontents
\vspace{1em}
\hrule
\vspace{1em}

% ── 1. Introduction ───────────────────────────────────────────────────────────
\section{Introduction}

Fine particulate matter (PM$_{2.5}$, diameter $\leq 2.5$~$\mu$m) is a leading
environmental health risk, associated with cardiovascular and respiratory disease.
Traditional monitoring networks rely on expensive regulatory-grade equipment with
sparse geographic coverage.
Low-cost optical particle counters, such as those used by PurpleAir, offer a
complementary approach: dense, community-scale networks at a fraction of the cost.

However, low-cost sensors introduce new challenges.
Individual units can drift, fail, or produce spurious readings without any
obvious external indication.
Raw sensor output must be cleaned, validated, and contextualized before it can
support public health decisions or research conclusions.

This project addresses those challenges directly.
We built a complete pipeline --- from raw CSV downloads to a deployed web dashboard ---
that (i)~scores the trustworthiness of every reading using the built-in dual-channel
redundancy of each sensor, (ii)~detects meaningful events such as smoke plumes,
precipitation washout, and sensor failures using domain-informed rules, and
(iii)~produces short-term PM$_{2.5}$ forecasts with calibrated uncertainty so that
users can see not just a prediction, but how confident the model is in it.

The system is designed to be understandable and reproducible.
All data are publicly available, all code is open source, and the dashboard can
be deployed to Vercel with a single command.

% ── 2. Related Work ───────────────────────────────────────────────────────────
\section{Background and Related Work}

\subsection{PurpleAir Sensors and Dual-Channel Architecture}

PurpleAir sensors embed two Plantower laser particle counters (channels A and B)
in each housing.
Both channels sample the same air simultaneously, providing built-in redundancy.
The EPA and PurpleAir recommend using channel agreement as a data quality indicator:
when A and B diverge significantly, the unit may be malfunctioning or contaminated.
Our trust score formalizes this intuition into a continuous 0--100 metric.

\subsection{EPA AQI and PM$_{2.5}$ Standards}

The U.S.\ Environmental Protection Agency updated its PM$_{2.5}$ AQI breakpoints
in February 2024, lowering the \textit{Good} category threshold from
12.0~$\mu$g/m$^3$ to 9.0~$\mu$g/m$^3$ to reflect the revised annual standard
\cite{epa2024}. Our system uses these updated 2024 breakpoints throughout.

\subsection{Time Series Forecasting for Air Quality}

Random Forest regressors with lag features have been shown to be competitive
with LSTM networks for short-horizon air quality forecasting, particularly when
training data are limited to a single season \cite{chen2021}.
We adopt this approach and augment it with meteorological features (temperature,
humidity) and temporal features (hour of day, day of week).

% ── 3. Dataset ────────────────────────────────────────────────────────────────
\section{Dataset}

\subsection{Data Source and Collection}

Data were collected by Dr.\ Kozminski's group at Lewis University and are
publicly available at \url{https://github.com/DrKoz/ACS_AQ}.
The dataset consists of five PurpleAir-style PM$_{2.5}$ sensor units deployed
during November 2025, recording at 2-minute intervals.

\subsection{Sensor Network}

Table~\ref{tab:sensors} summarizes the five sensors.
Sensor 290694 has substantially lower coverage (17.9\%), likely due to a
deployment delay or hardware issue; it is excluded from forecasting but
included in trust scoring and event detection.

\begin{table}[h!]
\centering
\caption{Sensor network summary --- November 1--29, 2025.}
\label{tab:sensors}
\begin{tabular}{lrr}
\toprule
\textbf{Sensor ID} & \textbf{Measurements} & \textbf{Coverage (\%)} \\
\midrule
230019 & 20,558 & 98.5 \\
249443 & 20,365 & 97.5 \\
249445 & 20,677 & 99.0 \\
290694 & 3,741  & 17.9 \\
297697 & 20,673 & 99.0 \\
\midrule
\textbf{Total} & \textbf{85,014} & --- \\
\bottomrule
\end{tabular}
\end{table}

\subsection{Variables}

Each row in the dataset contains the following fields:

\begin{itemize}[noitemsep]
  \item \textbf{\texttt{pm\_cf1}} --- EPA-corrected, averaged PM$_{2.5}$ ($\mu$g/m$^3$);
        primary target variable.
  \item \textbf{\texttt{corr\_pm2.5\_a}, \texttt{corr\_pm2.5\_b}} --- Corrected PM$_{2.5}$
        from channels A and B independently.
  \item \textbf{\texttt{temperature\_a}, \texttt{temperature\_b}} --- Ambient temperature
        from each channel.
  \item \textbf{\texttt{humidity\_a}, \texttt{humidity\_b}} --- Relative humidity from
        each channel.
  \item \textbf{\texttt{time\_stamp}} --- UTC timestamp at 2-minute resolution.
  \item \textbf{\texttt{sensor\_index}} --- Sensor identifier.
\end{itemize}

\subsection{Data Limitations}

The dataset covers a single calendar month, limiting seasonal generalizability.
Precipitation data and wind speed/direction are not included; both would
strengthen event detection and washout analysis.
Sensor GPS coordinates are not provided, which prevents spatial interpolation
or wind-direction-based smoke tracking.

% ── 4. Methods ────────────────────────────────────────────────────────────────
\section{Methods}

Our system consists of three independent components, each of which adds a layer
of information on top of the raw PM$_{2.5}$ readings.

\subsection{Component 1: Dual-Sensor Trust Scoring}

\subsubsection{Motivation}

Each sensor unit produces two independent PM$_{2.5}$ readings (channels A and B).
Under normal operation, the two channels should agree closely.
Persistent or large divergence indicates that at least one channel is unreliable,
whether due to sensor failure, contamination, or calibration drift.

\subsubsection{Formula}

For each timestamp, we compute a trust score $T \in [0, 100]$ as:

\begin{equation}
  T = 100 \cdot \exp\!\left(-3 \cdot \frac{|A - B|}{\bar{x}}\right),
  \quad \bar{x} = \frac{A + B}{2}
  \label{eq:trust}
\end{equation}

where $A = \mathtt{corr\_pm2.5\_a}$, $B = \mathtt{corr\_pm2.5\_b}$, and
$\bar{x}$ is their mean.
If $\bar{x} < 0.1$~$\mu$g/m$^3$ (both channels near zero), $T$ is set to 100
to avoid division-by-zero while respecting that near-zero agreement is valid.

\subsubsection{Interpretation}

The exponential decay ensures that small differences (e.g., 2\%) yield
near-perfect trust ($T \approx 94$), while large divergences (e.g., 67\%) collapse
trust toward zero.
Table~\ref{tab:trust_examples} illustrates the behavior.

\begin{table}[h!]
\centering
\caption{Trust score examples at various divergence levels.}
\label{tab:trust_examples}
\begin{tabular}{ccccc}
\toprule
$A$ ($\mu$g/m$^3$) & $B$ ($\mu$g/m$^3$) & $|A-B|/\bar{x}$ & $T$ & Badge \\
\midrule
10.0 & 10.2 & 2.0\%  & 94 & HIGH \\
10.0 & 12.0 & 18.2\% & 58 & MED  \\
10.0 & 20.0 & 66.7\% & 1  & LOW  \\
\bottomrule
\end{tabular}
\end{table}

\subsubsection{Results}

The mean trust score across all sensors and timestamps ranged from
78.2\% (sensor 290694) to 84.2\% (sensor 249445), consistent with
normally functioning dual-channel hardware.

% ─────────────────────────────────────────────────────────────────────────────

\subsection{Component 2: Rule-Based Event Detection}

\subsubsection{Motivation}

Three classes of meaningful events can be identified from the PM$_{2.5}$ time
series without labeled training data: (1)~\textit{smoke spikes} from fires or
pollution episodes, (2)~\textit{precipitation washout} events where rain rapidly
clears particulates, and (3)~\textit{sensor failures} where A/B divergence
indicates hardware problems.

\subsubsection{Preprocessing}

For each sensor, we compute a rolling mean $\mu_t$ and rolling change
$\Delta_t = x_t - x_{t-w}$ over a window $w$ (measured in 2-minute timesteps):
\[
  \mu_t = \frac{1}{w}\sum_{i=0}^{w-1} x_{t-i}, \qquad
  \Delta_t = x_t - x_{t-w}
\]

\subsubsection{Detection Rules}

\begin{itemize}[noitemsep]
  \item \textbf{Smoke spike:} $\Delta_t > 0.5 \cdot \mu_t$ and $x_t > 35$~$\mu$g/m$^3$
        (rapid rise $\geq$50\% above the 30-minute rolling mean, sustained above
        the EPA Unhealthy for Sensitive Groups threshold). Severity is \textit{high}
        if $x_t > 80$~$\mu$g/m$^3$, else \textit{medium}.
  \item \textbf{Rain washout:} $\Delta_t < -0.35 \cdot \mu_t$ and $\mu_t > 5$~$\mu$g/m$^3$
        (rapid drop $\geq$35\% over one hour from a non-trivial baseline).
  \item \textbf{Sensor failure:} $T < 40$ and $|A-B|/\bar{x} > 0.4$
        persisting for $\geq$3 consecutive readings ($\approx$6 min).
\end{itemize}

Events are deduplicated to at most one occurrence per type per sensor per
two-hour window to avoid alert flooding.

\subsubsection{Confidence Scoring}

Confidence is assigned heuristically:
\begin{align*}
  c_{\text{smoke}} &= \min\!\left(0.95,\; 0.60 + \frac{x_t - 35}{200}\right) \\
  c_{\text{washout}} &= \min\!\left(0.90,\; 0.50 + \frac{|\Delta_t|}{50}\right) \\
  c_{\text{failure}} &= 0.95
\end{align*}

These values reflect relative certainty: sensor failures that persist have
a high mechanical probability, while smoke and washout classifications may
be confounded by other sources (e.g., traffic, industrial emissions).

% ─────────────────────────────────────────────────────────────────────────────

\subsection{Component 3: PM$_{2.5}$ Forecasting with Uncertainty}

\subsubsection{Model}

We train a \textbf{Random Forest regressor} (100 trees, maximum depth 12,
scikit-learn default otherwise) per sensor to predict PM$_{2.5}$ at
timestep $t+1$ (two minutes ahead) from features available at time $t$.

\subsubsection{Feature Engineering}

The feature vector $\mathbf{x}_t$ is constructed entirely from past information,
with no future leakage:

\[
  \mathbf{x}_t = \left[
    x_{t-1},\, x_{t-2},\, \ldots,\, x_{t-12},\;
    \text{temp}_t,\;
    \text{humidity}_t,\;
    \text{hour}(t),\;
    \text{dow}(t)
  \right]
\]

where $x_{t-k}$ are 12 lagged PM$_{2.5}$ values (past 24 minutes),
$\text{temp}_t$ and $\text{humidity}_t$ are channel-averaged meteorological
readings, and $\text{hour}(t)$, $\text{dow}(t)$ capture diurnal and weekly
periodicity.

\subsubsection{Train/Test Split}

We use a strict \textbf{temporal split}: the first 80\% of each sensor's
time series is used for training, and the remaining 20\% is held out for
evaluation.
Random shuffling is never applied, preventing data leakage across the
temporal boundary.

\subsubsection{Uncertainty Estimation}

Prediction intervals are produced by \textbf{tree bootstrap}: each of
the 100 trees makes an independent prediction, yielding a distribution of
100 values per test point. The 80\% prediction interval is:

\[
  \left[\hat{x}_{0.10},\; \hat{x}_{0.90}\right]
\]

where $\hat{x}_{p}$ denotes the $p$-th percentile across trees.
This approach does not require a separate calibration set and is
computationally inexpensive given the model is already trained.

\subsubsection{Results}

Table~\ref{tab:forecast} reports forecast accuracy on the held-out test sets.

\begin{table}[h!]
\centering
\caption{PM$_{2.5}$ next-step forecast accuracy (test set, 20\% temporal holdout).}
\label{tab:forecast}
\begin{tabular}{lrrr}
\toprule
\textbf{Sensor} & \textbf{MAE ($\mu$g/m$^3$)} & \textbf{RMSE ($\mu$g/m$^3$)} & \textbf{$R^2$} \\
\midrule
230019 & 0.51 & 1.15 & 0.982 \\
249443 & 1.12 & 4.53 & 0.766 \\
249445 & 0.77 & 2.00 & 0.955 \\
297697 & 0.67 & 1.53 & 0.965 \\
\midrule
\textbf{Mean} & \textbf{0.77} & \textbf{2.30} & \textbf{0.917} \\
\bottomrule
\end{tabular}
\end{table}

Sensor 249443 shows higher error, likely due to greater variability in its
time series; its $R^2 = 0.77$ indicates that the model still explains the
majority of variance, but with less precision than the other units.

% ── 5. Dashboard ──────────────────────────────────────────────────────────────
\section{Dashboard}

\subsection{Architecture}

The monitoring dashboard is a React single-page application built with Vite.
A Python export script (\texttt{export\_dashboard\_data.py}) runs the full ML
pipeline and serializes the results to a static \texttt{data.json} file, which
the dashboard fetches on load.
This decoupled architecture means the computationally heavy Python pipeline
runs once at build time; the client-side dashboard is fast and requires no backend.
The application is deployable to Vercel with a single command.

\subsection{Tabs and Visualizations}

\begin{description}[style=nextline, leftmargin=2em]
  \item[\textbf{Overview}]
    Five sensor cards showing current AQI (EPA 2024 color scale), PM$_{2.5}$
    in $\mu$g/m$^3$, a trust badge (HIGH / MED / LOW), and a 60-point sparkline.
    A network aggregate panel shows the mean PM$_{2.5}$, maximum AQI, and
    count of recent alerts.

  \item[\textbf{Detail}]
    A full-width time series chart for a selected sensor.
    Time range is switchable (1h / 6h / 24h / 7d).
    EPA threshold lines (9, 35.4, 55.4~$\mu$g/m$^3$) are drawn as
    horizontal references.
    An optional A/B overlay shows channel-level readings alongside the
    fused \texttt{pm\_cf1} value, making sensor reliability directly visible.

  \item[\textbf{Events}]
    A chronological list of detected events with type (smoke spike /
    rain washout / sensor failure), sensor ID, timestamp, severity, and
    confidence score.
    An event-type summary bar shows aggregate counts.

  \item[\textbf{Forecast}]
    Per-sensor short-term forecast charts showing the Random Forest point
    prediction and 80\% tree-bootstrap confidence band.
    Sensor 290694 is flagged as unreliable and excluded from the network average.

  \item[\textbf{Map (Purple Air)}]
    An embedded iframe of the live PurpleAir regional PM$_{2.5}$ map,
    providing geographic context for the five-sensor network.
\end{description}

\subsection{EPA AQI Conversion}

AQI values are computed from \texttt{pm\_cf1} using the 2024 EPA piecewise
linear breakpoints \cite{epa2024}:

\begin{table}[h!]
\centering
\caption{EPA 2024 AQI breakpoints for PM$_{2.5}$ (24-hour average).}
\label{tab:aqi}
\begin{tabular}{llc}
\toprule
\textbf{AQI Range} & \textbf{Category} & \textbf{PM$_{2.5}$ ($\mu$g/m$^3$)} \\
\midrule
0--50   & Good                             & 0.0--9.0   \\
51--100 & Moderate                         & 9.1--35.4  \\
101--150 & Unhealthy for Sensitive Groups  & 35.5--55.4 \\
151--200 & Unhealthy                       & 55.5--125.4 \\
201--300 & Very Unhealthy                  & 125.5--225.4 \\
301--500 & Hazardous                       & 225.5+ \\
\bottomrule
\end{tabular}
\end{table}

% ── 6. Discussion ─────────────────────────────────────────────────────────────
\section{Discussion}

\subsection{Forecasting Performance}

The strong $R^2$ values (up to 0.98 for sensor 230019) confirm that next-step
PM$_{2.5}$ is highly predictable from recent history at 2-minute resolution.
This is expected: air quality changes slowly relative to the sampling interval,
so lag features carry most of the predictive signal.
The practical utility of next-step forecasting is limited; future work should
extend the horizon to 30 minutes or 1 hour and evaluate whether performance
degrades gracefully.

\subsection{Trust Score Limitations}

The trust formula in Equation~\ref{eq:trust} is analytically simple and requires
no training data, but it is not calibrated against an independent ground truth.
A reading could have high A/B agreement while both channels are simultaneously
biased (e.g., due to a contaminated housing).
Future work could calibrate the trust score against co-located reference monitors.

\subsection{Event Detection Without Labels}

The rule-based event detector identifies events by threshold logic, not by
learning from labeled examples.
Without a ground-truth event log (fire incident reports, precipitation gauge data),
it is not possible to report precision or recall.
The detector is best understood as a heuristic alert system that reduces the
search space for human analysts.

\subsection{Data Limitations}

The dataset covers a single month in a single location.
Seasonal patterns (e.g., winter heating, summer wildfires) and multi-year
trends are not captured.
The absence of precipitation data is a significant gap: without rain gauge
records, washout events cannot be definitively confirmed, only inferred from
PM$_{2.5}$ drops.

% ── 7. Conclusion ─────────────────────────────────────────────────────────────
\section{Conclusion}

We presented an end-to-end air quality monitoring system for a five-sensor
PM$_{2.5}$ network at Lewis University.
The system introduces a dual-sensor trust scoring metric based on channel
agreement, a rule-based event detection pipeline for smoke, washout, and
sensor failure events, and a Random Forest forecasting model with tree-bootstrap
uncertainty intervals.
Forecast accuracy ranges from MAE 0.51--1.12~$\mu$g/m$^3$ ($R^2$ 0.77--0.98)
across the four high-coverage sensors.
All components are integrated into a React dashboard deployed on Vercel,
with real data exported from the Python pipeline at build time.

Future directions include extending the forecast horizon, calibrating trust
scores against reference monitors, incorporating external precipitation and
wind data, and exploring multi-sensor fusion to estimate network-wide PM$_{2.5}$
fields.

% ── Acknowledgements ──────────────────────────────────────────────────────────
\section*{Acknowledgements}

The authors thank Dr.\ Kozminski at Lewis University for providing sensor data,
project guidance, and continued mentorship throughout this work.

% ── References ────────────────────────────────────────────────────────────────
\section*{References}
\begin{thebibliography}{9}

\bibitem{epa2024}
U.S. Environmental Protection Agency.
\textit{Final Updates to the Air Quality Index (AQI) for Particulate Matter ---
Fact Sheet and Common Questions}.
EPA-456/F-24-001. February 2024.
\url{https://www.epa.gov/system/files/documents/2024-02/pm-naaqs-air-quality-index-fact-sheet.pdf}

\bibitem{chen2021}
Chen, T., \& Guestrin, C.
\textit{XGBoost: A scalable tree boosting system}.
Proceedings of the 22nd ACM SIGKDD International Conference on Knowledge
Discovery and Data Mining, 2016, pp.\ 785--794.

\bibitem{purpleair}
PurpleAir, Inc.
\textit{PurpleAir Real-Time Air Quality Map}.
\url{https://map.purpleair.com/}, accessed 2026.

\bibitem{acs_data}
Kozminski, D.
\textit{ACS\_AQ: Air Quality Sensor Network Data --- November 2025}.
GitHub repository.
\url{https://github.com/DrKoz/ACS_AQ}, 2025.

\bibitem{scikit}
Pedregosa, F., et al.
\textit{Scikit-learn: Machine learning in Python}.
Journal of Machine Learning Research, 12, pp.\ 2825--2830, 2011.

\end{thebibliography}

\end{document}
